\documentclass[11pt]{article}
\usepackage[a4paper,margin=1in]{geometry}
\usepackage{amsmath,amssymb,amsthm,mathtools}
\usepackage{bm}
\usepackage{booktabs}
\usepackage{enumitem}
\usepackage[dvipsnames]{xcolor}
\usepackage{hyperref}
\usepackage{microtype}
\usepackage{comment}

\hypersetup{
  colorlinks=true,
  linkcolor=MidnightBlue,
  urlcolor=MidnightBlue
}

% Define \WITHSOLUTIONS before inputting this file to show solutions.
\newif\ifsolutions
\ifdefined\WITHSOLUTIONS
  \solutionstrue
\else
  \solutionsfalse
\fi

\newcommand{\Pbb}{\mathbb{P}}
\newcommand{\E}{\mathbb{E}}
\newcommand{\Var}{\operatorname{Var}}
\newcommand{\iid}{\mathrel{\overset{\mathrm{iid}}{\sim}}}
\newcommand{\MLE}{\mathrm{ML}}
\newcommand{\R}{\mathbb{R}}
\newcommand{\rv}[1]{\widetilde{#1}}

\newtheorem*{definition}{Definition}

\ifsolutions
  \newenvironment{answer}{%
    \par\medskip\noindent\textcolor{MidnightBlue}{\textbf{Solution.}}\ }
    {\par\medskip}
  \newcommand{\workingspace}[1]{}
\else
  \excludecomment{answer}
  \newcommand{\workingspace}[1]{\vspace{#1}}
\fi

\title{\huge Recitation 1\\[0.35em]
  \normalsize Maximum Likelihood Estimation\\[0.35em]
  \large\ifsolutions Solutions and Teaching Notes\else Review and Practice Problems\fi}
\author{Xuanxi Zhang}
\date{}

\begin{document}
\maketitle

\ifsolutions
\section*{Pacing guide (75 minutes)}
\begin{center}
\begin{tabular}{@{}lll@{}}
\toprule
Time & Activity & Suggested use \\
\midrule
0--25 min & MLE review & Setup, Poisson derivation, Gaussian derivation \\
25--37 min & Problem 1 & 7 min individual work, then discussion \\
37--50 min & Problem 2 & 7 min individual work, then discussion \\
50--66 min & Problem 3 & 8 min individual work, then discussion \\
66--75 min & Wrap-up or Problem 4 & Use Problem 4 only if the class is ahead \\
\bottomrule
\end{tabular}
\end{center}
\fi

\section{Maximum likelihood estimation: the basic setup}

\paragraph{Course notation.}
A tilde marks a random quantity: $\rv{x}_1,\ldots,\rv{x}_n$. After observing
these random variables, their realizations are $x_1,\ldots,x_n$, and
$X:=\{x_1,\ldots,x_n\}$ denotes the dataset. Thus, in this course, $X$ is a
dataset rather than a random variable. Model parameters and values fitted to
the observed dataset are written without tildes or hats.

Suppose $\rv{x}_1,\ldots,\rv{x}_n$ are i.i.d. random variables from a model
indexed by a parameter $\theta\in\Theta$, and let
$X=\{x_1,\ldots,x_n\}$ be the observed dataset. The \textbf{likelihood} is
\[
  L(\theta;X)
  = \prod_{i=1}^n p_\theta(x_i)
  \quad\text{or}\quad
  L(\theta;X)
  = \prod_{i=1}^n f_\theta(x_i),
\]
depending on whether the model is discrete or continuous. Once the data are
observed, $L$ is a function of $\theta$; it is not a pmf or pdf in $\theta$.

\begin{definition}
The maximum likelihood (ML) estimate for the observed dataset $X$ is
\[
  \theta_{\MLE}
  \in \underset{\theta\in\Theta}{\arg\max}\,L(\theta;X).
\]
Because $\log$ is strictly increasing, maximizing $L$ is equivalent to
maximizing the log-likelihood
\[
  \ell(\theta)=\log L(\theta)=\sum_{i=1}^n \log p_\theta(x_i)
  \quad\text{or}\quad
  \ell(\theta)=\sum_{i=1}^n \log f_\theta(x_i).
\]
\end{definition}

\noindent A reliable workflow is:
\begin{enumerate}[label=\arabic*.,leftmargin=2em,itemsep=0.2em]
  \item state the parameter space and the model assumptions;
  \item write the likelihood, including any parameter-dependent support;
  \item take logs and remove only terms that do not depend on the parameter;
  \item find candidate maximizers and check the boundary or support;
  \item report the estimate in the units of the original parameter.
\end{enumerate}

\section{Poisson MLE}

Let $\rv{x}_1,\ldots,\rv{x}_n\iid\operatorname{Poisson}(\lambda)$, where
$\lambda\geq 0$. For the observed counts
$X=\{x_1,\ldots,x_n\}$,
\begin{align*}
  L(\lambda)
  &=\prod_{i=1}^n e^{-\lambda}\frac{\lambda^{x_i}}{x_i!},\\
  \ell(\lambda)
  &=-n\lambda+\left(\sum_{i=1}^n x_i\right)\log\lambda
    -\sum_{i=1}^n\log(x_i!).
\end{align*}
If $\sum_i x_i>0$, then
\[
  \ell'(\lambda)=-n+\frac{\sum_i x_i}{\lambda},
  \qquad
  \ell''(\lambda)=-\frac{\sum_i x_i}{\lambda^2}<0.
\]
Therefore
\[
  \boxed{\lambda_{\MLE}=\frac{1}{n}\sum_{i=1}^n x_i=m(X)=\overline x.}
\]
The estimate has a natural interpretation: the fitted expected count equals
the sample average. If every observed count is zero, the likelihood decreases
as $\lambda$ increases, so the maximizer is the boundary value
$\lambda_{\MLE}=0$.

\section{Gaussian MLE}

Let $\rv{x}_1,\ldots,\rv{x}_n\iid\mathcal N(\mu,\sigma^2)$, where
$\mu\in\R$ and $\sigma^2>0$, and let $X=\{x_1,\ldots,x_n\}$ be the observed
dataset. Write $v=\sigma^2$. The log-likelihood is
\[
  \ell(\mu,v)
  =-\frac n2\log(2\pi)-\frac n2\log v
    -\frac{1}{2v}\sum_{i=1}^n(x_i-\mu)^2.
\]
For fixed $v$,
\[
  \frac{\partial\ell}{\partial\mu}
  =\frac{1}{v}\sum_{i=1}^n(x_i-\mu)=0
  \quad\Longrightarrow\quad
  \boxed{\mu_{\MLE}=m(X)=\overline x.}
\]
The identity
\[
  \sum_{i=1}^n(x_i-\mu)^2
  =\sum_{i=1}^n(x_i-\overline x)^2+n(\mu-\overline x)^2
\]
also shows directly that $\mu=\overline x$ gives the global maximum for fixed
$v$. Substitute $\mu_{\MLE}=\overline x$ and let
$S=\sum_i(x_i-\overline x)^2$. Then
\[
  \ell(\overline x,v)=C-\frac n2\log v-\frac{S}{2v},
  \qquad
  \frac{d\ell}{dv}=-\frac{n}{2v}+\frac{S}{2v^2}.
\]
Thus, when $S>0$,
\[
  \boxed{\sigma_{\MLE}^2
  =\frac{1}{n}\sum_{i=1}^n(x_i-\overline x)^2,}
  \qquad
  \boxed{\sigma_{\MLE}
  =\sqrt{\frac{1}{n}\sum_{i=1}^n(x_i-\overline x)^2}.}
\]

\paragraph{MLE variance versus the usual sample variance.}
The Gaussian MLE divides by $n$. The unbiased sample variance divides by
$n-1$. Maximum likelihood targets the parameter value that best fits the
observed data; it does not guarantee an unbiased estimator.

\paragraph{Degenerate data.}
If all observations are equal, then $S=0$. With parameter space $v>0$, the
likelihood grows without bound as $v\downarrow0$, so no MLE for $v$ exists in
that open parameter space. This is a useful reminder to check the boundary.

\newpage
\section*{Practice problems}

\subsection*{Problem 1: A two-point model \hfill [7 minutes]}
Suppose $\rv{x}$ is a discrete random variable with
\[
  \Pbb(\rv{x}=1)=\theta,
  \qquad
  \Pbb(\rv{x}=2)=1-\theta,
  \qquad 0\leq\theta\leq1.
\]
We draw $n$ independent observations and find that $n_1$ of them equal 1 and
$n_2$ equal 2, where $n_1+n_2=n$.
\begin{enumerate}[label=(\alph*),itemsep=0.35em]
  \item Write the likelihood function.
  \item Find the maximum likelihood estimate of $\theta$, including the
        cases $n_1=0$ or $n_2=0$.
  \item Rewrite the MLE in terms of the sample mean $\overline x$.
\end{enumerate}
\workingspace{4.5cm}

\begin{answer}
The likelihood is
\[
  L(\theta)=\theta^{n_1}(1-\theta)^{n_2}.
\]
For $n_1,n_2>0$,
\[
  \ell'(\theta)=\frac{n_1}{\theta}-\frac{n_2}{1-\theta}=0
  \quad\Longrightarrow\quad
  \theta_{\MLE}=\frac{n_1}{n}.
\]
Also $\ell''(\theta)<0$ on $(0,1)$. The same formula covers the boundaries:
$n_1=0$ gives $\theta_{\MLE}=0$, and $n_2=0$ gives $\theta_{\MLE}=1$.
Finally,
\[
  \overline x=\frac{n_1+2n_2}{n}=2-\frac{n_1}{n}
  \quad\Longrightarrow\quad
  \theta_{\MLE}=2-\overline x.
\]
\end{answer}

\subsection*{Problem 2: Calls to a help desk \hfill [7 minutes]}
The numbers of calls received in eight comparable 15-minute periods are
\[
  3,\ 0,\ 2,\ 1,\ 4,\ 2,\ 0,\ 4.
\]
Assume the counts are realizations of
$\rv{x}_1,\ldots,\rv{x}_8\iid\operatorname{Poisson}(\lambda)$.
\begin{enumerate}[label=(\alph*),itemsep=0.35em]
  \item Find the MLE of $\lambda$.
  \item Using the fitted model, estimate the probability that a future
        15-minute period has at least three calls.
  \item Give one modeling assumption that would make the calculation in
        part (b) questionable.
\end{enumerate}
\workingspace{4.5cm}

\begin{answer}
The total count is 16, so
\[
  \lambda_{\MLE}=\overline x=\frac{16}{8}=2
  \quad\text{calls per 15 minutes}.
\]
If $\rv{y}\sim\operatorname{Poisson}(2)$ is the number of calls in a future
period, then
\[
  \begin{aligned}
  \Pbb_{\lambda_{\MLE}}(\rv{y}\geq3)
  &=1-\Pbb_{\lambda_{\MLE}}(\rv{y}=0)
    -\Pbb_{\lambda_{\MLE}}(\rv{y}=1)
    -\Pbb_{\lambda_{\MLE}}(\rv{y}=2)\\
  &=1-e^{-2}\left(1+2+\frac{2^2}{2}\right)
   =1-5e^{-2}\approx0.323.
  \end{aligned}
\]
Possible concerns include dependence across adjacent periods, different call
rates at different times, or overdispersion relative to a Poisson model. Any
one well-explained concern is sufficient.
\end{answer}

\newpage
\subsection*{Problem 3: Repeated measurements \hfill [8 minutes]}
A device records the same reference quantity five times:
\[
  9.7,\ 10.1,\ 10.2,\ 9.9,\ 10.1.
\]
Assume these values are realizations of
$\rv{x}_1,\ldots,\rv{x}_5\iid\mathcal N(\mu,\sigma^2)$.
\begin{enumerate}[label=(\alph*),itemsep=0.35em]
  \item Compute $\mu_{\MLE}$.
  \item Compute $\sigma_{\MLE}^2$ and $\sigma_{\MLE}$.
  \item Compute the usual unbiased sample variance $s^2$. Why does its
        denominator differ from the MLE denominator?
  \item If 2 is added to every observation, how do the two MLEs change?
\end{enumerate}
\workingspace{6.0cm}

\begin{answer}
The sample mean is
\[
  \mu_{\MLE}=\overline x=10.0.
\]
The centered observations are $-0.3,0.1,0.2,-0.1,0.1$, whose squared values
sum to $0.16$. Therefore
\[
  \sigma_{\MLE}^2=\frac{0.16}{5}=0.032,
  \qquad
  \sigma_{\MLE}=\sqrt{0.032}\approx0.179.
\]
The unbiased sample variance is
\[
  s^2=\frac{0.16}{4}=0.04.
\]
Estimating $\mu$ uses one degree of freedom, so division by $n-1$ corrects the
downward bias. Maximum likelihood instead maximizes the Gaussian likelihood
and therefore divides by $n$. Adding 2 to every observation changes
$\mu_{\MLE}$ from $10.0$ to $12.0$ and leaves $\sigma_{\MLE}^2$ unchanged.
\end{answer}

\ifsolutions\else\newpage\fi
\subsection*{Problem 4 (backup): True or false?}
\noindent\textit{Optional if time allows. Suggested working time: 5--8 minutes.}

\noindent For each statement, decide whether it is true or false and give a brief reason.
\par\smallskip
\begin{enumerate}[label=(\alph*),itemsep=0.55em]
  \item MLEs are always unbiased.
  \item If $\rv{\theta}_n\xrightarrow{p}\theta$, then
        $\rv{\theta}_n^2\xrightarrow{p}\theta^2$.
  \item Every solution of the score equation $\ell'(\theta)=0$ is an MLE.
  \item A likelihood must satisfy
        $\int_{-\infty}^{\infty}L(\theta)\,d\theta=1$.
\end{enumerate}
\workingspace{4.5cm}

\begin{answer}
The answers are false, true, false, false.
\begin{enumerate}[label=(\alph*),itemsep=0.35em]
  \item The Gaussian MLE of $\sigma^2$ is a counterexample: its expectation is
        $\frac{n-1}{n}\sigma^2$.
  \item This follows from the continuous mapping theorem because
        $g(t)=t^2$ is continuous. The premise explicitly assumes consistency.
  \item A zero derivative only identifies a stationary point. It may be a
        local minimum, and it can miss a maximizer on the boundary.
  \item A likelihood is a function of the parameter with the data held fixed.
        It need not integrate to one over the parameter space. A normalized
        density in $\theta$ arises only after adding a prior and forming a
        posterior in a Bayesian analysis.
\end{enumerate}
\end{answer}

\section*{Takeaways}
\begin{itemize}[leftmargin=1.6em,itemsep=0.25em]
  \item For i.i.d. data, likelihoods multiply and log-likelihoods add.
  \item The Poisson MLE matches the sample mean count.
  \item The Gaussian MLE matches the sample mean and uses divisor $n$ for the
        variance.
  \item A stationary point is not enough: check the parameter space, boundary,
        and parameter-dependent support.
\end{itemize}

\section*{Sources}
The review follows the course handouts
\href{https://github.com/cfgranda/ps4ds/blob/main/slides/discrete%20variables/parametric_vs_nonparametric_handout.pdf}{Comparing Parametric and Nonparametric Models}
and
\href{https://github.com/cfgranda/ps4ds/blob/main/slides/continuous%20variables/maximum_likelihood_continuous_models_handout.pdf}{Maximum Likelihood Estimation for Continuous Models}.
The notation follows Section 2.1 of the
\href{https://github.com/cfgranda/ps4ds/blob/main/ps4ds_preprint.pdf}{Probability and Statistics for Data Science preprint}.
Problem 1 and Problem 4 are adapted from the 2025 Fall statistics recitation
materials (Recitation 5).

\end{document}
